\begin{abstract}
激光超声检测技术可实现非接触式探伤，但往往面临信噪比（SNR）较低的问题，需要进行大量信号平均处理，这不仅降低了检测通量，还存在表面损伤的风险。现有去噪方法难以充分保留对缺陷检测至关重要的声学特征，限制了其实际应用价值。本文提出PMDF-Net，一种基于物理约束的多域融合网络，通过联合编码时域波形和时频表示来恢复高质量超声信号。该框架采用双分支架构，并行处理时域波形和时频谱图，并借助跨域融合模块自适应地整合多尺度特征。物理约束损失函数通过可微分惩罚函数强制执行到达时间、峰值幅值和波形形态的一致性，确保保留声学相关的信号分量。训练采用配对数据，包括低平均次数输入（N=16平均）和高平均次数目标（M=256平均），使网络能够从有限测量中学习信号的本质特征。在激光超声基准数据集上的实验表明，PMDF-Net实现了11.4 dB的信噪比提升，同时在特征保留质量上表现最优（CCC=0.981），并超越了基线方法的缺陷检测F1分数。这使得平均次数需求减少了16倍，同时保留了与高平均参考方法相当的缺陷检测性能，有效缩短了检测时间并减少了表面损伤。此外，本文还验证了从完好样本到含缺陷工件的零样本泛化能力。本文方法在保持缺陷检测能力的前提下，显著降低了激光超声检测的时间和表面损伤风险。
\end{abstract}